% Canadian AI 2027 - position paper (6 pages)
% FORMAT NOTE: Canadian AI has not published its 2027 CFP. This file uses
% IEEEtran conference format as a working layout. When the CFP appears,
% switch to whatever it requires (Springer LNCS is the likely one).
% DOUBLE-BLIND NOTE: the 2026 edition used double-blind review. Before
% submitting, comment out the \author block below and anonymise ref [25].
\documentclass[conference]{IEEEtran}
\usepackage[T1]{fontenc}
\usepackage[utf8]{inputenc}
\usepackage{amsmath,amssymb}
\usepackage{array}
\usepackage{url}
\usepackage[hidelinks]{hyperref}
\usepackage{microtype}
\usepackage{balance}
\begin{document}
\title{Grounding the AI Scrum Master: Evidence Before Autonomy}
\author{\IEEEauthorblockN{Hani AbuSharkh}
\IEEEauthorblockA{Higher Colleges of Technology\\Al Ain, UAE\\habusharkh@hct.ac.ae}}
\date{}
\maketitle

\begin{abstract}
Agent systems carrying the Scrum Master label are multiplying, but few of them do Scrum Master work, and none that we found has been tested against what actually happened in real sprints. We place the systems we found against the accountabilities the role carries. The evidence clusters on estimation and planning; facilitation, impediment removal, review and retrospective, and coaching are close to empty. A separate line of work predicts sprint outcomes from real Jira histories, but stops at prediction and never acts. Our position is that these two lines should meet, and that the meeting point sets a testable bar: an agent in this role should ground what it says in models validated on real sprint data, and should be judged by replaying historical sprints without hindsight against a model-only arm and an agent-only arm. We set out what that evaluation requires, and argue that two failure modes bracket the design space. An ungrounded agent states things the data does not support. An over-grounded one repeats its model and adds nothing. Neither would survive the bar we propose, which is why the bar is worth setting now, while the label is still ahead of the evidence.
\end{abstract}

\begin{IEEEkeywords}
LLM agents, Scrum, agile software development, sprint analytics, evaluation, grounding
\end{IEEEkeywords}

\section{Introduction}

Much of what a Scrum Master does never reaches the tracker. A grounded-theory study of the role names six everyday functions: facilitating, mentoring, negotiating, adapting processes, coordinating, and protecting the team \cite{shastri}. Only two of those leave much of a trace in Jira. The rest happen in conversations, in standing between a team and an interruption, in noticing that a retrospective has gone quiet.

Language-model agents have nonetheless started to wear the title. Within the last two years, systems have been published that assign an agent the Scrum Master role in a code-generation pipeline \cite{agilecoder}, simulate agile teams with role agents \cite{cognisim}, run stand-ups for student teams \cite{couder}, and turn burndown charts into status reports \cite{alliata}. Tool vendors have gone further: agents can now be assigned issues directly in commercial trackers, and at least one company reports an in-house AI Scrum Master that runs rituals by voice.

Our position is that the label has run ahead of the evidence, and that it can be brought back into line with a specific, affordable test. The test is sprint replay: rebuild a project's board as it stood at a chosen moment, give the agent only what was knowable then, let it decide what to say, and score that against what the sprint went on to do. Nothing in this requires a partner organisation or a live team. What it does require is a sensing layer the agent can cite, and a willingness to report how often the agent merely repeats it.

We argue three things. First, that the systems wearing the label mostly do not do the job, and that for the parts of the job a tracker can see the evidence is thin while for the rest it is close to absent (Sections~\ref{sec:label} and \ref{sec:gap}). Second, that two documented failure modes bracket what a credible system has to avoid (Section~\ref{sec:failure}). Third, that grounding in leakage-controlled models, judged by hindsight-free replay, is a sufficient and affordable bar, and we say what the replay has to include to count (Sections~\ref{sec:grounding} and \ref{sec:eval}).

\section{The label and the job}
\label{sec:label}

The systems that use Scrum roles fall into three families, and the distance from the role grows as the family changes.

The first family organises code generation. AgileCoder gives a Scrum Master agent the job of reviewing a backlog, choosing sprint tasks and judging deliverability, and is scored on code benchmarks alone \cite{agilecoder}. A system from an industry lab merges Product Manager and Scrum Master into a single sprint agent that splits tasks and estimates story points, evaluated on one small application and one feature request \cite{almas}. In both, the Scrum Master does no facilitation, no impediment work and no coaching. The role supplies a vocabulary, not a set of responsibilities.

The second family simulates agile teams. CogniSim models scaled-agile roles with language-model agents and, notably, includes no Scrum Master agent at all; its own future work asks for predictive analytics for delays \cite{cognisim}.

The third family targets the role directly, and here the evidence is real but small. An action-research study at a large postal organisation put assistants into daily Scrums and refinement for three teams; about a third of early refinement outputs were wrong because the input data was stale, two of seven people in one team found the assistant annoying, and participants asked that it inform rather than warn \cite{cabrero}. Another study prompts models to produce status reports from burndown charts and scores the text with similarity metrics; there is no agent and no prediction \cite{alliata}. A third uses a model as stand-up facilitator for student capstone teams \cite{couder}.

Industry is ahead on plumbing and behind on evidence. Integration with trackers is becoming a commodity, and the built-in agents mostly do housekeeping: splitting epics, tidying backlogs, writing meeting notes. The open question is not whether an agent can act in a tracker. It is what it should say, when, and whether that helps.

\section{Where the evidence is, and is not}
\label{sec:gap}

Recent reviews agree on the shape of the field. A review of 395 papers on language models in software engineering puts software management at under one percent of usages, against more than half for development \cite{hou}. A mapping study finds planning dominates its use cases, with only a handful of studies on agile meetings \cite{planotscher}. A review of 18 studies finds no applications at all in Scrum's implementation and review-and-retrospect phases \cite{peras}. A fourth notes that most work treats the model as a low-autonomy tool and calls for it to become an active team member \cite{kwok}.

Table~\ref{tab:gap} places the role's accountabilities against the best evidence we found for each.

\begin{table}[!t]
\centering\scriptsize
\caption{Scrum Master accountabilities against available evidence. ``Real data'' means real project records used offline; ``field'' means people in real teams.}
\label{tab:gap}
\setlength{\tabcolsep}{3pt}
\begin{tabular}{llp{2.1cm}}
\hline
Accountability & Coverage & Best evidence\\
\hline
Effort / story-point estimation & Well covered & Real data, many projects \cite{tse2018,sprint2vec,tawosi23}\\
Backlog refinement & Moderate & Real stories; small field \cite{cabrero}\\
Delivery-risk prediction at planning & Moderate & Real Jira \cite{tse2018,sprint2vec,helwerda}\\
Sprint planning / scope selection & Thin & One proprietary optimiser \cite{kula24}\\
In-sprint monitoring & Thin & Checkpoints, no agent \cite{tse2018,kula23}\\
Daily Scrum facilitation & Thin & One field site \cite{cabrero}; classroom \cite{couder}\\
Impediment detection & Thin & Link detection, no role framing \cite{re2022}\\
Impediment removal & Empty & None academic\\
Review / stakeholder reporting & Near empty & Text scoring \cite{alliata}; simulation \cite{agentcompany}\\
Retrospective facilitation & Near empty & Small, author-labelled \cite{enase}\\
Coaching, team health & Empty & Detectors only \cite{toxicr}\\
End-to-end agent in the role & Label only & Benchmarks and demos \cite{agilecoder,almas}\\
\hline
\end{tabular}
\end{table}

Three patterns follow. First, the more relational the task, the thinner the evidence; of the six functions the role actually performs \cite{shastri}, only facilitation and coordination have any at all. Second, evaluation habits are weak: small samples, single runs, author-labelled ground truth. Run-to-run instability recurs wherever anyone measures it \cite{yoon}. Third, the data is lopsided. Public sprint histories are plentiful, with one widely used dataset holding thousands of sprints and millions of change-log entries \cite{tawos}; public transcripts of stand-ups, reviews and retrospectives are, as far as we can tell, absent. Any proposal to automate the relational half of this role runs into that absence immediately. The half that is left, the half the tracker sees, is where an agent can be held to account today.

\section{Two ways this fails}
\label{sec:failure}

Two failure modes bracket the design space, each with a published result behind it, and a system in this role has to sit between them.

\subsection{Ungrounded}

Asked Scrum knowledge questions, current models do well on recall and badly on interpretation; in one evaluation over several hundred certification-style items, at least one prompt failed on roughly a sixth of them, with errors clustering on interpretive topics \cite{psm}. Asked to turn natural language into tracker queries, one assistant was correct between 17 and 49 percent of the time and invented project names until the prompt was constrained \cite{jql}. Asked to produce sprint plans, models were judged not of acceptable quality for direct use, with priority accuracy swinging from 44 to 92 percent across runs \cite{yoon}. In a mission-critical risk setting, retrieval-grounded models had the lowest hallucination rates \cite{esem24}. The lesson is not that the models are unusable. It is that in this role, confident wrong statements are the characteristic failure, and grounding is the characteristic fix.

\subsection{Over-grounded}

The opposite failure is less discussed and, for research purposes, more dangerous. When agents are given a predictive model as a tool, they tend to copy it. In one study, agents reproduced a graph model's output on between 97.6 and 99.2 percent of cases, and larger models deferred more rather than less \cite{tool}. An agent that repeats its model is not an agent contribution; it is the model with extra latency. Any paper claiming an agent in this role should therefore report its agreement rate with the raw model, and should include a model-only arm so that readers can see what the agent added. We would go further: a system that agrees with its own tool on nearly every case has not demonstrated judgement, whatever its task scores say. It has demonstrated the tool.

None of this is particular to Scrum. Agents are now being placed over classical predictive models in clinical triage, credit decisions and operations, and the same two questions arise in each: when should the agent believe the model, and when is acting on it worth the cost? Sprints are an unusually good place to study both, because the ground truth is public, time-stamped, and free of hindsight once the board is rebuilt. Whatever is learned here about deference and timing should travel.

\section{What grounding would mean here}
\label{sec:grounding}

Grounding, in this setting, means the agent's statements about risk trace back to models fitted on real sprint histories under evaluation conditions that do not leak the future.

That last clause carries most of the weight. Sprint-outcome studies report accuracies and areas under the curve from 0.87 to 0.9995, often using features derived from the outcome itself, or applying resampling before splitting the data. A stricter pipeline produces lower numbers, and the lower numbers are the trustworthy ones. Leakage taxonomies now list software engineering among the affected fields \cite{leakage}. An agent that cites an inflated model inherits the inflation and adds a confident voice to it.

Grounding also constrains what the agent may suggest. Work on defect prediction found that guidance derived from explanation methods often does not fix the thing it predicts, and differs from what developers actually change \cite{saner}. Methods that restrict proposed changes to those with historical precedent reach far higher overlap with real developer behaviour \cite{counteract}. Carried into sprints, this means an agent should not propose a scope reduction simply because the model's expectation falls; it should propose changes that teams in comparable situations actually made.

What remains for the agent, then, is narrower than the label suggests and more interesting than the model alone: reading the issue text and comments the tabular model never sees, deciding whether the situation is worth raising at all, choosing the moment, and wording it so that a team can act without feeling audited. Each of those is measurable.

\section{How to evaluate it}
\label{sec:eval}

We propose sprint replay as the primary instrument, because it is the only evaluation we can run at scale today without a partner organisation.

\textbf{Construction.} Public datasets carry change logs of field-level edits with timestamps \cite{tawos}, which is enough to rebuild a board as it stood at any point in a sprint. Prior work has reconstructed such snapshots for prediction \cite{tse2018}. The agent sees only events recorded at or before the replay point.

\textbf{Arms.} At minimum three: the risk model alone, the agent with no model, and the grounded agent. Without the first, no one can tell whether the agent added anything; without the second, no one can tell whether the grounding did.

\textbf{Measures.} Detection quality against realised outcomes; agreement with the raw model, reported as a first-class result \cite{tool}; calibration, since a ranked warning and a probability are different products; and cost-sensitive timing. The last is borrowed from prescriptive process monitoring, where alarm policies trade intervention cost against failure cost and the value of a warning decays as a case proceeds \cite{alarm}. A warning at a quarter of the way through a sprint is worth more than the same warning on the final day, and both cost something when wrong.

\textbf{Reliability.} Repeated runs with variance reported. Single-run results in this literature are common and are not evidence.

\textbf{What replay cannot do.} It measures what an agent would have said, not whether saying it would have changed anything. The counterfactual is not in the log. Establishing effect requires people, and the one controlled result we know of on language-model facilitation found that it surfaced more information, mostly by drawing in quieter participants, yet did not improve decision quality \cite{alsobay}. Replay is a filter, not a verdict: a system that fails it should not reach a team, and a system that passes it has earned a study, not a deployment.

\section{Limits of this position}

Public sprint data is open-source, with volunteer staffing and uneven administrative discipline; enterprise teams may behave differently. Completion timestamps record when someone closed a sprint, not when work finished, so outcomes built on them are administrative events rather than performance measures. Models may have seen public tracker data during pre-training, which argues for an agent-only arm and for replication outside software data. And the relational half of the role stays out of reach while no public ceremony corpus exists. We regard that as a reason to be explicit about scope rather than a reason to wait: the half a tracker records is enough to settle whether grounding helps, and that question is in front of us now.

Finally, autonomy. Developers surveyed on how much autonomy they would give AI keep planning and management work at low levels \cite{choudhuri}, and the one field study in this role reports participants asking to be informed rather than warned, with one asking that the Scrum Master approve what the assistant shares \cite{cabrero}. An agent designed to inform and prepare is both more likely to be accepted and easier to evaluate honestly than one designed to decide.

\balance

\section{Conclusion}

The AI Scrum Master exists as a label, a vendor feature and a handful of pilots. It does not yet exist as a system whose judgement has been measured against what really happened in sprints. Meanwhile a decade of sprint-prediction work sits on the other side of the same gap, sensing without acting. Our position is that the two belong together, and that the join can be tested cheaply: replay real sprints, ground the agent in models whose evaluation does not leak the future, and report how often the agent simply repeated them. The bar is low enough to clear with public data and high enough that most current systems would not.

\begin{thebibliography}{33}
\bibitem{shastri} Y. Shastri, R. Hoda, and R. Amor, ``Spearheading agile: the role of the scrum master in agile projects,'' Empir. Softw. Eng., vol. 26, no. 1, Art. no. 3, 2021.
\bibitem{agilecoder} M. H. Nguyen, T. P. Chau, P. X. Nguyen, and N. D. Q. Bui, ``AgileCoder: dynamic collaborative agents for software development based on agile methodology,'' in Proc. IEEE/ACM Int. Workshop AI Found. Softw. Eng. (FORGE), 2025.
\bibitem{cognisim} J. Cinkusz and J. A. Chudziak, ``Agile software management with cognitive multi-agent systems,'' in Proc. 17th Int. Conf. Agents Artif. Intell. (ICAART), 2025.
\bibitem{couder} J. Ortiz Couder and O. Ochoa, ``Large language models, the new scrum masters,'' in Proc. EDULEARN24, 2024, pp. 9928--9934.
\bibitem{alliata} G. Alliata, A. Singhal, and R. Bozagiu, ``The AI scrum master,'' in Agile Processes in Software Engineering and Extreme Programming -- Workshops (XP 2024), LNBIP 524, Springer, 2024.
\bibitem{almas} V. Tawosi et al., ``ALMAS: an autonomous LLM-based multi-agent software engineering framework,'' arXiv:2510.03463, 2025. (preprint)
\bibitem{cabrero} B. Cabrero-Daniel et al., ``Exploring human-AI collaboration in agile: customised LLM meeting assistants,'' in Agile Processes in Software Engineering and Extreme Programming (XP 2024), LNBIP, Springer, 2024.
\bibitem{hou} X. Hou et al., ``Large language models for software engineering: a systematic literature review,'' ACM Trans. Softw. Eng. Methodol., vol. 33, no. 8, Art. no. 220, 2024.
\bibitem{planotscher} M. Planötscher, ``Large language models in agile software development: a systematic mapping study,'' in Agile Processes in Software Engineering and Extreme Programming -- Workshops, Springer, 2026.
\bibitem{peras} D. Peras et al., ``Artificial intelligence in Scrum: a systematic literature review,'' Informatica, 2026.
\bibitem{kwok} A. Kwok and M. Adil, ``From tool to team member: a review of AI in agile software development,'' in Agile Processes in Software Engineering and Extreme Programming -- Workshops, Springer, 2026.
\bibitem{tse2018} M. Choetkiertikul, H. K. Dam, T. Tran, A. Ghose, and J. Grundy, ``Predicting delivery capability in iterative software development,'' IEEE Trans. Softw. Eng., vol. 44, no. 6, pp. 551--573, Jun. 2018.
\bibitem{sprint2vec} M. Choetkiertikul, P. Banyongrakkul, C. Ragkhitwetsagul, S. Tuarob, H. K. Dam, and T. Sunetnanta, ``Sprint2Vec: a deep characterization of sprints in iterative software development,'' IEEE Trans. Softw. Eng., vol. 51, no. 1, pp. 220--242, Jan. 2025.
\bibitem{tawosi23} V. Tawosi, R. Moussa, and F. Sarro, ``Agile effort estimation: have we solved the problem yet? Insights from a replication study,'' IEEE Trans. Softw. Eng., vol. 49, no. 4, pp. 2677--2697, Apr. 2023.
\bibitem{helwerda} L. Helwerda, F. Niessink, and F. Verbeek, ``Grip on software: understanding development progress of SCRUM sprints and backlogs,'' in Proc. 9th Int. Conf. Knowl. Discovery Inf. Retrieval (KDIR), 2017.
\bibitem{kula24} E. Kula, E. Greuter, A. van Deursen, and G. Gousios, ``Context-aware automated sprint plan generation for agile software development,'' in Proc. 39th IEEE/ACM Int. Conf. Automated Softw. Eng. (ASE), 2024.
\bibitem{kula23} E. Kula, A. van Deursen, and G. Gousios, ``Dynamic prediction of delays in software projects using delay patterns and Bayesian modeling,'' in Proc. 31st ACM Joint Eur. Softw. Eng. Conf. Symp. Found. Softw. Eng. (ESEC/FSE), 2023.
\bibitem{tawos} V. Tawosi, A. Al-Subaihin, R. Moussa, and F. Sarro, ``A versatile dataset of agile open source software projects,'' in Proc. 19th Int. Conf. Mining Softw. Repositories (MSR), 2022, pp. 707--711.
\bibitem{re2022} C. M. Lüders, T. Pietz, and W. Maalej, ``Automated detection of typed links in issue trackers,'' in Proc. IEEE 30th Int. Requirements Eng. Conf. (RE), 2022.
\bibitem{agentcompany} F. F. Xu et al., ``TheAgentCompany: benchmarking LLM agents on consequential real world tasks,'' in Adv. Neural Inf. Process. Syst. (NeurIPS), 2025.
\bibitem{enase} M. Spichkova, H. Lee, K. Iwan, M. Zwart, Y. Yoon, and X. Qin, ``Agile retrospectives: what went well? What didn't go well? What should we do?,'' in Proc. 20th Int. Conf. Eval. Novel Approaches Softw. Eng. (ENASE), 2025, pp. 746--753.
\bibitem{toxicr} J. Sarker, A. K. Turzo, M. Dong, and A. Bosu, ``Automated identification of toxic code reviews using ToxiCR,'' ACM Trans. Softw. Eng. Methodol., vol. 32, no. 5, Art. no. 118, 2023.
\bibitem{psm} M. Perkusich et al., ``Evaluating large language models on professional scrum master assessments,'' arXiv:2607.00049, 2026. (preprint)
\bibitem{jql} J. Garcia-Escribano, A. Carbajo, M. Egaña Aranguren, and U. Lopez-Novoa, ``Using GPT to build a project management assistant for Jira environments,'' arXiv:2509.26014, 2025. (preprint)
\bibitem{yoon} Y. Yoon, K. Iwan, M. Zwart, X. Qin, H. Lee, and M. Spichkova, ``Scrum sprint planning: LLM-based and algorithmic solutions,'' arXiv:2512.18966, 2025. (preprint)
\bibitem{esem24} M. Esposito, F. Palagiano, V. Lenarduzzi, and D. Taibi, ``Beyond words: on large language models actionability in mission-critical risk analysis,'' in Proc. 18th ACM/IEEE Int. Symp. Empirical Softw. Eng. Meas. (ESEM), 2024.
\bibitem{tool} Z. Wang and P. Vemuri, ``When the tool decides: LLM agents defer blindly to graph neural network tools, and stronger backbones defer more,'' arXiv:2606.14476, 2026. (preprint)
\bibitem{leakage} S. Kapoor and A. Narayanan, ``Leakage and the reproducibility crisis in machine-learning-based science,'' Patterns, vol. 4, no. 9, Art. no. 100804, 2023.
\bibitem{saner} G. Lee and S. U.-J. Lee, ``Can we trust the actionable guidance from explainable AI techniques in defect prediction?,'' in Proc. IEEE Int. Conf. Softw. Anal., Evolution Reengineering (SANER), 2025.
\bibitem{counteract} K. Oueslati, G. Laberge, M. Lamothe, and F. Khomh, ``Mining action rules for defect reduction planning,'' in Proc. ACM Int. Conf. Found. Softw. Eng. (FSE), 2024.
\bibitem{alarm} S. A. Fahrenkrog-Petersen, N. Tax, I. Teinemaa, M. Dumas, M. de Leoni, F. M. Maggi, and M. Weidlich, ``Fire now, fire later: alarm-based systems for prescriptive process monitoring,'' Knowl. Inf. Syst., vol. 64, no. 2, pp. 559--587, 2022.
\bibitem{alsobay} M. Alsobay, D. M. Rothschild, J. M. Hofman, and D. G. Goldstein, ``Bringing everyone to the table: an experimental study of LLM-facilitated group decision making,'' arXiv:2508.08242, 2025. (preprint)
\bibitem{choudhuri} R. Choudhuri, C. Bird, C. Badea, M. Gerosa, and A. Sarma, ``You shall not pass! Where and why developers draw the line on AI autonomy,'' arXiv:2607.00533, 2026. (preprint)
\end{thebibliography}

\end{document}
